% from ega3-6-fr.tex
\oldpage[III]{142}

et nous dirons que cet $A$-module est l'hypertor d'indice $n$ des deux
complexes $P_\bullet$, $Q_\bullet$. On sait que dans la catégorie des
complexes de $A$-modules limités inférieurement, les
$\mathbf{Tor}_n^A(P_\bullet,Q_\bullet)$ forment un bifoncteur homologique en
$P_\bullet$, $Q_\bullet$ (0, 11.6.5). En outre :
\end{env}

\begin{proposition}[6.3.2]
Le bifoncteur $\mathbf{Tor}_\bullet^A(P_\bullet,Q_\bullet)$ est
l'aboutissement commun de deux bifoncteurs spectraux
$'E(P_\bullet,Q_\bullet)$, $''E(P_\bullet,Q_\bullet)$, dont les termes $E_2$
sont
\begin{align}
  'E^2_{pq}
    &=H_p\bigl(\operatorname{Tor}_q^A(P_\bullet,Q_\bullet)\bigr),
    \tag{6.3.2.1}\label{III.6.3.2.1.fr}\\
  ''E^2_{pq}
    &=\bigoplus_{q'+q''=q}
      \operatorname{Tor}_p^A\bigl(H_{q'}(P_\bullet),H_{q''}(Q_\bullet)\bigr).
    \tag{6.3.2.2}\label{III.6.3.2.2.fr}
\end{align}
où, dans (6.3.2.1), $\operatorname{Tor}_q^A(P_\bullet,Q_\bullet)$ désigne le
bicomplexe formé des $A$-modules $\operatorname{Tor}_q^A(P_i,Q_j)$. La suite
spectrale (6.3.2.2) est toujours régulière; si $P_\bullet$ et $Q_\bullet$
sont limités inférieurement, ou si $A$ est de dimension cohomologique finie,
les deux suites spectrales (6.3.2.1) et (6.3.2.2) sont birégulières.
\end{proposition}

\begin{proof}
Cela résulte de (0, 11.6.5), car lorsque $A$ est de dimension cohomologique
finie $n$, tout $A$-module admet une résolution projective de longueur $n$
(M, VI, 2.1).
\end{proof}

\begin{corollary}[6.3.3]
Soient $P'_\bullet$, $Q'_\bullet$ deux complexes de $A$-modules,
$u:P_\bullet\to P'_\bullet$, $v:Q_\bullet\to Q'_\bullet$ deux homomorphismes
de complexes. Si les homomorphismes
$H_\bullet(u):H_\bullet(P_\bullet)\to H_\bullet(P'_\bullet)$,
$H_\bullet(v):H_\bullet(Q_\bullet)\to H_\bullet(Q'_\bullet)$ déduits
respectivement de $u$ et $v$ sont bijectifs, alors l'homomorphisme
$\mathbf{Tor}_\bullet^A(P_\bullet,Q_\bullet)\to
\mathbf{Tor}_\bullet^A(P'_\bullet,Q'_\bullet)$ déduit de $u$ et $v$ est
bijectif.
\end{corollary}

\begin{proof}
En effet, l'homomorphisme de suites spectrales
$''E(P_\bullet,Q_\bullet)\to''E(P'_\bullet,Q'_\bullet)$ déduit de $u$ et $v$
est alors un isomorphisme pour les termes $E_2$ et la conclusion résulte de
ce que ces suites sont régulières en vertu de (6.3.2) (0, 11.1.5).
\end{proof}

\begin{proposition}[6.3.4]
Soient $P_\bullet$, $Q_\bullet$ deux complexes de $A$-modules, limités
inférieurement. Soit $L_{\bullet\bullet}$ (resp. $M_{\bullet\bullet}$) un
bicomplexe formé de $A$-modules plats, tel que pour tout $i$,
$L_{i,\bullet}$ (resp. $M_{i,\bullet}$) soit une résolution de $P_i$
(resp. $Q_i$). On a alors des isomorphismes canoniques
\[
  \mathbf{Tor}_\bullet^A(P_\bullet,Q_\bullet)
  \xrightarrow{\sim}H_\bullet(L_{\bullet\bullet}\otimes_AQ_\bullet)
  \xrightarrow{\sim}H_\bullet(P_\bullet\otimes_AM_{\bullet\bullet})
  \xrightarrow{\sim}H_\bullet(L_{\bullet\bullet}\otimes_AM_{\bullet\bullet})
  \tag{6.3.4.1}\label{III.6.3.4.1.fr}
\]
\end{proposition}

\begin{proof}
Cela résulte de (0, 11.6.5, (ii) et (iii)) et de la définition des
$A$-modules plats.
\end{proof}

\begin{namedtheorem}[6.3.5]{Remarques}
(i) Avec les notations de (6.3.1), les bicomplexes
$L_{\bullet\bullet}\otimes_A M_{\bullet\bullet}$ et
$M_{\bullet\bullet}\otimes_A L_{\bullet\bullet}$ sont canoniquement
isomorphes, d'où un isomorphisme canonique
$\mathbf{Tor}_\bullet^A(P_\bullet,Q_\bullet)\xrightarrow{\sim}
\mathbf{Tor}_\bullet^A(Q_\bullet,P_\bullet)$.

(ii) Si $F$ et $G$ sont deux $A$-modules, $P_\bullet$ et $Q_\bullet$ les
complexes de $A$-modules réduits à $F$ et $G$ respectivement en degré 0 et
nuls dans les autres degrés, alors deux résolutions projectives $L_\bullet$,
$M_\bullet$ de $F$ et $G$ respectivement peuvent être considérées comme des
résolutions de Cartan-Eilenberg de $P_\bullet$ et $Q_\bullet$ en les
complétant par des zéros. On a par suite dans ce cas
$\mathbf{Tor}_\bullet^A(P_\bullet,Q_\bullet)
=\operatorname{Tor}_\bullet^A(F,G)$.
\end{namedtheorem}

\begin{proposition}[6.3.6]
Soient $(P_\bullet^\lambda)$, $(Q_\bullet^\mu)$ deux systèmes inductifs
filtrants de complexes de $A$-modules; on a un isomorphisme canonique
\[
  \varinjlim_{\lambda,\mu}
  \mathbf{Tor}_\bullet^A(P_\bullet^\lambda,Q_\bullet^\mu)
  \xrightarrow{\sim}
  \mathbf{Tor}_\bullet^A
  \left(\varinjlim_\lambda P_\bullet^\lambda,
        \varinjlim_\mu Q_\bullet^\mu\right).
  \tag{6.3.6.1}\label{III.6.3.6.1.fr}
\]
\end{proposition}

\begin{proof}
Posons $P_\bullet=\varinjlim_\lambda P_\bullet^\lambda$,
$Q_\bullet=\varinjlim_\mu Q_\bullet^\mu$; par fonctorialité, il est clair que
les $\mathbf{Tor}_\bullet^A(P_\bullet^\lambda,Q_\bullet^\mu)$ forment un
système inductif et que les applications
$\mathbf{Tor}_\bullet^A(P_\bullet^\lambda,Q_\bullet^\mu)\to
\mathbf{Tor}_\bullet^A(P_\bullet,Q_\bullet)$ déduites

